\documentclass[12pt, a4paper]{article}
\usepackage[utf8]{inputenc}
\usepackage[T1]{fontenc}
\usepackage[french]{babel}
\usepackage{amsmath, amssymb, amsthm}
\usepackage{geometry}
\geometry{hmargin=2.5cm, vmargin=2.5cm}
\usepackage{listings}
\usepackage{tikz}
\usepackage{xcolor}
\usepackage{hyperref}

\lstset{
    language=Python,
    basicstyle=\ttfamily\small,
    keywordstyle=\color{blue},
    stringstyle=\color{red},
    commentstyle=\color{green!60!black},
    showstringspaces=false,
    breaklines=true,
    frame=single,
    backgroundcolor=\color{gray!10}
}

\title{Formes linéaires, hyperplans, espace dual et orthogonalité en dimension finie}
\author{Charles EDOU NZE}
\date{\today}

\begin{document}

\maketitle
\tableofcontents
\newpage


\section{Jalon 11 : Formes linéaires, hyperplans, espace dual et
orthogonalité en dimension
finie}


\subsection{1. Présentation du concept
clé}

La genèse des formes linéaires trouve ses racines dans le besoin
inhérent à l'être humain de mesurer, quantifier et extraire une
information scalaire (unidimensionnelle) à partir d'objets
multidimensionnels complexes. Imaginez une vaste plaine topographique,
où chaque position est déterminée par une multitude de coordonnées
complexes. Un explorateur cherchant à quantifier l'altitude de son
terrain de progression effectue sans le savoir une projection de ce
monde multidimensionnel vers l'axe des réels : il applique une forme
linéaire.

Historiquement, cette idée d'associer un unique scalaire à un vecteur
s'est cristallisée lorsque les mathématiciens ont cherché à comprendre
la structure intime des espaces eux-mêmes, non pas à travers les
vecteurs qui les composent, mais à travers le prisme de toutes les
``mesures'' possibles sur ces vecteurs. C'est l'essence de la dualité,
une notion formellement structurée au début du XXe siècle par des
figures comme Stefan Banach et David Hilbert. Si l'on considère un
espace vectoriel \(E\) comme un univers de points, l'espace dual \(E^*\)
représente l'univers de tous les instruments de mesure, ou observateurs,
capables de scruter \(E\). Un hyperplan, concept fondamental de
séparation, émerge naturellement : il s'agit de l'horizon, la frontière
parfaitement plate (de dimension \(n-1\)) où une mesure spécifique
s'annule, séparant l'espace en deux demi-espaces stricts.


\subsection{2. Formalisation}

L'intuition de la mesure laisse place à une rigueur algébrique
implacable. Nous posons ici les fondations algébriques de la dualité en
dimension finie.


\subsubsection{A. Anatomie et Typage
Chirurgical}

Soit \(E\) un espace vectoriel sur un corps commutatif \(\mathbb{K}\)
(typiquement \(\mathbb{R}\) ou \(\mathbb{C}\)), de dimension finie
\(\dim(E) = n \in \mathbb{N}^*\).

\begin{enumerate}

\item
  \textbf{Forme linéaire :} Une forme linéaire est une application
  \(\phi : E \to \mathbb{K}\) satisfaisant la propriété de linéarité
  stricte :
  \[ \forall (x, y) \in E^2, \forall (\lambda, \mu) \in \mathbb{K}^2, \quad \phi(\lambda x + \mu y) = \lambda \phi(x) + \mu \phi(y) \]
  Le typage est fondamental : la source est l'espace vectoriel \(E\), et
  le but est le corps de base \(\mathbb{K}\) (considéré lui-même comme
  un \(\mathbb{K}\)-espace vectoriel de dimension 1).
\item
  \textbf{Espace Dual (\(E^*\)) :} L'ensemble de toutes les formes
  linéaires sur \(E\), noté \(E^* = \mathcal{L}(E, \mathbb{K})\), est
  l'espace dual de \(E\). Il hérite d'une structure naturelle de
  \(\mathbb{K}\)-espace vectoriel.
\item
  \textbf{Hyperplan :} Un sous-espace vectoriel \(H \subset E\) est
  appelé un hyperplan si et seulement s'il existe une forme linéaire
  \textbf{non nulle} \(\phi \in E^* \setminus \{0_{E^*}\}\) telle que :
  \[ H = \ker(\phi) = \{ x \in E \mid \phi(x) = 0_{\mathbb{K}} \} \]
  L'exigence \(\phi \neq 0_{E^*}\) est impérative pour exclure le cas
  trivial où \(H = E\).
\item
  \textbf{Base Duale :} Si l'on fixe une base
  \(\mathcal{B} = (e_1, e_2, \dots, e_n)\) de l'espace vectoriel \(E\),
  on définit l'unique base duale
  \(\mathcal{B}^* = (e_1^*, e_2^*, \dots, e_n^*)\) de \(E^*\) par la
  relation de dualité de Kronecker :
  \[ \forall (i, j) \in \llbracket 1, n \rrbracket^2, \quad e_i^*(e_j) = \delta_{i,j} = \begin{cases} 1_{\mathbb{K}} & \text{si } i = j \\ 0_{\mathbb{K}} & \text{si } i \neq j \end{cases} \]
\item
  \textbf{Orthogonalité Duale :} Soit \(A\) une partie non vide de
  \(E\). Son orthogonal dans l'espace dual \(E^*\) est le sous-espace :
  \[ A^\perp = \{ \phi \in E^* \mid \forall x \in A, \phi(x) = 0_{\mathbb{K}} \} \]
  Réciproquement, si \(B \subset E^*\), son orthogonal prédual dans
  \(E\) est
  \(B^\circ = \{ x \in E \mid \forall \phi \in B, \phi(x) = 0_{\mathbb{K}} \}\).
\end{enumerate}


\subsubsection{B. Théorèmes
Fondamentaux}

\textbf{Théorème 1 (Dimension de l'espace dual) :} Si \(E\) est de
dimension finie, alors \(E\) et son dual \(E^*\) ont la même dimension :
\[ \dim(E) < +\infty \implies \dim(E^*) = \dim(E) \]

\textbf{Théorème 2 (Isomorphisme Canonique au Bidual) :} Le bidual est
défini comme \(E^{**} = (E^*)^*\). L'application d'évaluation
\(J : E \to E^{**}\) définie par
\(\forall x \in E, J(x)(\phi) = \phi(x)\) (souvent noté
\(\langle \phi, x \rangle\)) est un isomorphisme canonique. Cet
isomorphisme est naturel car il ne dépend d'aucun choix de base,
contrairement aux isomorphismes entre \(E\) et \(E^*\).


\subsection{3. Démonstrations}

Nous détaillons ici intégralement la preuve de l'isomorphisme canonique
entre l'espace \(E\) et son bidual \(E^{**}\), pierre angulaire de
l'algèbre linéaire en dimension finie.

\textbf{Théorème :} Soit \(E\) un \(\mathbb{K}\)-espace vectoriel de
dimension finie \(n\). L'application \(J : E \to E^{**}\) définie par
\(\forall x \in E, \forall \phi \in E^*, J(x)(\phi) = \phi(x)\) est un
isomorphisme de \(\mathbb{K}\)-espaces vectoriels.

\textbf{Preuve pas-à-pas :}

\begin{enumerate}

\item
  \textbf{Typage et bonne définition de \(J\) :} Fixons \(x \in E\).
  L'application \(J(x)\) prend en argument une forme linéaire
  \(\phi \in E^*\) et renvoie un scalaire \(\phi(x) \in \mathbb{K}\).
  Vérifions que \(J(x)\) est linéaire (c'est-à-dire
  \(J(x) \in E^{**}\)). Soient \(\phi, \psi \in E^*\) et
  \(\lambda \in \mathbb{K}\) :
  \[ J(x)(\lambda \phi + \psi) = (\lambda \phi + \psi)(x) \] Par
  définition des opérations sur \(\mathcal{L}(E, \mathbb{K})\) :
  \[ (\lambda \phi + \psi)(x) = \lambda \phi(x) + \psi(x) = \lambda J(x)(\phi) + J(x)(\psi) \]
  Donc \(J(x)\) est bien une forme linéaire sur \(E^*\). L'application
  globale \(J\) est donc bien typée : \(J : E \to E^{**}\).
\item
  \textbf{Linéarité de \(J\) :} Montrons que \(J\) est une application
  linéaire. Soient \(x, y \in E\) et \(\alpha \in \mathbb{K}\). Il faut
  montrer l'égalité fonctionnelle dans \(E^{**}\) :
  \(J(\alpha x + y) = \alpha J(x) + J(y)\). Pour toute forme
  \(\phi \in E^*\), évaluons les deux membres :
  \[ J(\alpha x + y)(\phi) = \phi(\alpha x + y) \] Par linéarité de
  \(\phi\) (car \(\phi \in E^*\)) :
  \[ \phi(\alpha x + y) = \alpha \phi(x) + \phi(y) \] D'autre part, la
  définition des opérations dans \(E^{**}\) donne :
  \[ (\alpha J(x) + J(y))(\phi) = \alpha J(x)(\phi) + J(y)(\phi) = \alpha \phi(x) + \phi(y) \]
  L'égalité sur chaque \(\phi \in E^*\) prouve que
  \(J(\alpha x + y) = \alpha J(x) + J(y)\). L'application \(J\) est donc
  linéaire.
\item
  \textbf{Injectivité de \(J\) :} Étudions le noyau \(\ker(J)\). Soit
  \(x \in \ker(J)\). Cela signifie que \(J(x) = 0_{E^{**}}\),
  c'est-à-dire que pour toute
  \(\phi \in E^*, J(x)(\phi) = 0_{\mathbb{K}}\). Donc,
  \(\forall \phi \in E^*, \phi(x) = 0_{\mathbb{K}}\). Supposons par
  l'absurde que \(x \neq 0_E\). Puisque \(x\) est un vecteur non nul
  d'un espace de dimension finie, le théorème de la base incomplète
  autorise à étendre le vecteur unique \((x)\) en une base de \(E\) :
  \(\mathcal{B} = (x, e_2, \dots, e_n)\). Considérons alors la première
  forme linéaire coordonnée \(e_1^*\) de la base duale
  \(\mathcal{B}^*\). Par construction, \(e_1^*(x) = 1_{\mathbb{K}}\).
  Or, nous avions déduit que pour toute forme \(\phi\),
  \(\phi(x) = 0_{\mathbb{K}}\). En appliquant cela à \(\phi = e_1^*\),
  on obtient \(1_{\mathbb{K}} = 0_{\mathbb{K}}\), ce qui est une
  contradiction manifeste dans un corps. L'hypothèse \(x \neq 0_E\) est
  donc fausse. On en déduit que \(x = 0_E\), d'où \(\ker(J) = \{0_E\}\).
  \(J\) est injective.
\item
  \textbf{Bijectivité par argument dimensionnel :} Nous savons que si
  \(\dim(E) = n\), alors \(\dim(E^*) = n\). En appliquant ce même
  théorème à l'espace vectoriel \(E^*\) (qui est aussi de dimension
  finie \(n\)), on obtient
  \(\dim(E^{**}) = \dim((E^*)^*) = \dim(E^*) = n\). Ainsi,
  \(\dim(E) = \dim(E^{**}) = n\). Puisque \(J : E \to E^{**}\) est une
  application linéaire injective entre deux espaces de même dimension
  finie, c'est obligatoirement un isomorphisme. La démonstration est
  achevée. \(\blacksquare\)
\end{enumerate}


\subsection{4. Exercices d'Application}

\emph{(Les exercices et corrections exhaustives sont répartis dans le
répertoire \texttt{exos/}.)}


\subsection{5. Application en Intelligence
Artificielle}

La dualité n'est pas qu'une abstraction algébrique ; c'est le langage
géométrique fondamental de la classification supervisée en apprentissage
automatique. Dans l'algorithme des Séparateurs à Vaste Marge (Support
Vector Machines - SVM), la séparation de données linéairement séparables
dans \(\mathbb{R}^n\) se formule précisément comme la recherche d'un
hyperplan optimal.

Un hyperplan de décision affine est défini par une équation de la forme
\(\phi(x) + b = 0\), où \(x \in \mathbb{R}^n\) est le vecteur de
caractéristiques, \(\phi \in (\mathbb{R}^n)^*\) est une forme linéaire
(qui, par le théorème de représentation de Riesz s'identifie au produit
scalaire avec un vecteur normal \(w\),
\(\phi(x) = \langle w, x \rangle\)), et \(b \in \mathbb{R}\) un biais.

La phase d'apprentissage d'un SVM consiste à trouver la forme linéaire
\(\phi\) qui maximise la marge, c'est-à-dire la distance entre
l'hyperplan \(\ker(\phi) - b\) et les points d'entraînement les plus
proches (les vecteurs de support). Le théorème d'isomorphisme dual est
exploité lors de la transformation du problème d'optimisation primal
complexe (minimisation de \(\|w\|^2\) sous contraintes) vers sa
formulation duale de Lagrange. C'est dans cet espace dual que le fameux
``Kernel Trick'' (astuce du noyau) opère, permettant d'évaluer
indirectement des formes linéaires dans des espaces de Hilbert de
dimension infinie sans jamais calculer explicitement les coordonnées des
vecteurs dans ces espaces.


\subsection{6. Liens Sémantiques}

\begin{itemize}

\item
  \textbf{Concepts Précédents requis :} {[}{[}Jalon-7.md\textbar Jalon 7
  (Espaces vectoriels abstraits){]}{]}, {[}{[}Jalon-8.md\textbar Jalon 8
  (Applications linéaires){]}{]}
\item
  \textbf{Concepts Futurs dépendants :} {[}{[}Jalon 12 (Livrable
  IA).md\textbar Jalon 12 (Livrable IA){]}{]},
  {[}{[}Jalon-25.md\textbar Jalon 25 (Formes bilinéaires){]}{]},
  {[}{[}Jalon-123.md\textbar Jalon 123 (Problèmes d'optimisation sous
  contraintes){]}{]}
\end{itemize}


\section{Exercices d'Application}


\section{Exercice 1: Espace dual en dimension
2}


\subsection{Énoncé}

Soit \(E = \mathbb{R}^2\). On considère la base canonique
\((e_1, e_2)\). Déterminer la base duale \((e_1^*, e_2^*)\) et calculer
\(e_1^*(3e_1 - 2e_2)\).


\subsection{Correction détaillée}

\begin{enumerate}


\item
  \textbf{Définition de la base duale :} Par définition, la base duale
  \((e_1^*, e_2^*)\) d'une base \((e_1, e_2)\) vérifie
  \(e_i^*(e_j) = \delta_{ij}\).
\item
  \textbf{Explicitation des formes linéaires :}

  \begin{itemize}

  \item
    \(e_1^*(x, y) = e_1^*(xe_1 + ye_2) = x e_1^*(e_1) + y e_1^*(e_2) = x \cdot 1 + y \cdot 0 = x\).
  \item
    \(e_2^*(x, y) = e_2^*(xe_1 + ye_2) = x e_2^*(e_1) + y e_2^*(e_2) = x \cdot 0 + y \cdot 1 = y\).
  \end{itemize}
\item
  \textbf{Calcul de l'évaluation :} On cherche à évaluer \(e_1^*\) sur
  le vecteur \(v = 3e_1 - 2e_2\).
\item
  \textbf{Développement complet :}
  \[e_1^*(3e_1 - 2e_2) = 3 e_1^*(e_1) - 2 e_1^*(e_2)\]
  \[= 3 \times 1 - 2 \times 0\] \[= 3\]
\item
  \textbf{Conclusion :} La valeur de la forme linéaire \(e_1^*\) sur le
  vecteur \(3e_1 - 2e_2\) est \(3\).
\end{enumerate}

\(\blacksquare\) \# Exercice 2: Noyau d'une forme linéaire \#\# Énoncé
Soit \(\phi : \mathbb{R}^3 \to \mathbb{R}\) définie par
\(\phi(x, y, z) = 2x - y + 3z\). Déterminer une base de \(\ker \phi\).


\subsection{Correction détaillée}

\begin{enumerate}


\item
  \textbf{Définition du noyau :} Le noyau de \(\phi\) est l'ensemble des
  vecteurs sur lesquels la forme linéaire s'annule.
  \[\ker \phi = \{ (x,y,z) \in \mathbb{R}^3 \mid 2x - y + 3z = 0 \}\]
\item
  \textbf{Résolution de l'équation cartésienne :} L'équation
  \(2x - y + 3z = 0\) équivaut à exprimer une variable en fonction des
  autres. Isolons \(y\) : \[y = 2x + 3z\]
\item
  \textbf{Paramétrisation des vecteurs du noyau :} Un vecteur
  \(v \in \ker \phi\) s'écrit donc : \[v = (x, 2x+3z, z)\]
\item
  \textbf{Décomposition en combinaison linéaire :} On sépare les
  paramètres libres \(x\) et \(z\) :
  \[v = (x, 2x, 0) + (0, 3z, z) = x(1, 2, 0) + z(0, 3, 1)\]
\item
  \textbf{Famille génératrice :} Les vecteurs \(u_1 = (1, 2, 0)\) et
  \(u_2 = (0, 3, 1)\) engendrent donc \(\ker \phi\).
\item
  \textbf{Indépendance linéaire :} Supposons
  \(\lambda u_1 + \mu u_2 = 0\).
  \[\lambda(1, 2, 0) + \mu(0, 3, 1) = (\lambda, 2\lambda+3\mu, \mu) = (0, 0, 0)\]
  On obtient immédiatement \(\lambda = 0\) et \(\mu = 0\). La famille
  \((u_1, u_2)\) est libre.
\item
  \textbf{Conclusion :} La famille \(((1,2,0), (0,3,1))\) est une base
  de \(\ker \phi\), qui est bien un hyperplan (dimension \(3 - 1 = 2\)).
\end{enumerate}

\(\blacksquare\) \# Exercice 3: Orthogonal d'un sous-espace \#\# Énoncé
Soit \(F\) un sous-espace vectoriel d'un \(\mathbb{K}\)-espace vectoriel
\(E\) de dimension finie \(n\). Démontrer le théorème de la dimension de
l'orthogonal : \(\dim F + \dim F^\perp = n\).


\subsection{Correction détaillée}

\begin{enumerate}


\item
  \textbf{Définition de l'application restriction :} On considère
  l'application linéaire \(\rho : E^* \to F^*\) définie par
  \(\rho(\phi) = \phi_{|F}\), c'est-à-dire que pour toute forme linéaire
  \(\phi\) sur \(E\), \(\rho(\phi)\) est sa restriction au sous-espace
  \(F\).
\item
  \textbf{Étude du noyau de \(\rho\) :} Par définition,
  \[\ker \rho = \{ \phi \in E^* \mid \rho(\phi) = 0 \} = \{ \phi \in E^* \mid \forall x \in F, \phi(x) = 0 \}\]
  Or, l'ensemble des formes linéaires qui s'annulent sur \(F\) est
  exactement la définition de l'orthogonal \(F^\perp\). Donc
  \(\ker \rho = F^\perp\).
\item
  \textbf{Étude de l'image de \(\rho\) :} L'application \(\rho\) est
  surjective. En effet, soit \(\psi \in F^*\). On peut compléter une
  base \((e_1, \dots, e_p)\) de \(F\) en une base
  \((e_1, \dots, e_p, e_{p+1}, \dots, e_n)\) de \(E\). On définit alors
  une forme \(\phi \in E^*\) par \(\phi(e_i) = \psi(e_i)\) pour
  \(1 \le i \le p\) et \(\phi(e_i) = 0\) pour \(i > p\). Ainsi
  \(\phi_{|F} = \psi\). Donc \(\text{Im}(\rho) = F^*\).
\item
  \textbf{Application du théorème du rang :} Le théorème du rang
  appliqué à \(\rho : E^* \to F^*\) donne :
  \[\dim E^* = \dim \ker \rho + \dim \text{Im}(\rho)\]
\item
  \textbf{Substitutions des dimensions :}

  \begin{itemize}

  \item
    \(\dim E^* = \dim E = n\)
  \item
    \(\dim \ker \rho = \dim F^\perp\)
  \item
    \(\dim \text{Im}(\rho) = \dim F^* = \dim F\) On obtient :
    \(n = \dim F^\perp + \dim F\).
  \end{itemize}
\item
  \textbf{Conclusion :} L'égalité \(\dim F + \dim F^\perp = n\) est
  rigoureusement démontrée.
\end{enumerate}

\(\blacksquare\) \# Exercice 4: Hyperplans et formes proportionnelles
\#\# Énoncé Montrer que deux formes linéaires non nulles \(\phi\) et
\(\psi\) ont le même noyau si et seulement si elles sont
proportionnelles.


\subsection{Correction détaillée}

\begin{enumerate}


\item
  \textbf{Sens direct (proportionnalité implique même noyau) :}
  Supposons qu'il existe un scalaire \(\lambda \neq 0\) tel que
  \(\phi = \lambda \psi\). Soit \(x \in \ker \phi\). Alors
  \(\phi(x) = 0\), donc \(\lambda \psi(x) = 0\). Comme
  \(\lambda \neq 0\), on a \(\psi(x) = 0\), donc \(x \in \ker \psi\).
  Ainsi \(\ker \phi \subset \ker \psi\). Symétriquement,
  \(\psi = \frac{1}{\lambda} \phi\), ce qui donne
  \(\ker \psi \subset \ker \phi\). Donc \(\ker \phi = \ker \psi\).
\item
  \textbf{Sens réciproque (même noyau implique proportionnalité) :}
  Supposons que \(\ker \phi = \ker \psi = H\). Comme \(\phi \neq 0\), il
  existe un vecteur \(e_0 \in E\) tel que \(\phi(e_0) \neq 0\). Quitte à
  diviser par \(\phi(e_0)\), on peut choisir \(e_0\) tel que
  \(\phi(e_0) = 1\).
\item
  \textbf{Décomposition d'un vecteur :} Soit \(x \in E\) quelconque.
  Posons \(x' = x - \phi(x)e_0\).
\item
  \textbf{Évaluation de \(x'\) :} Évaluons \(\phi\) en \(x'\) :
  \[\phi(x') = \phi(x - \phi(x)e_0) = \phi(x) - \phi(x)\phi(e_0) = \phi(x) - \phi(x)(1) = 0\]
  Donc \(x' \in \ker \phi\).
\item
  \textbf{Utilisation de l'égalité des noyaux :} Puisque
  \(\ker \phi = \ker \psi\), on a nécessairement \(x' \in \ker \psi\),
  ce qui implique \(\psi(x') = 0\).
\item
  \textbf{Relation de proportionnalité :}
  \[\psi(x - \phi(x)e_0) = 0 \implies \psi(x) - \phi(x)\psi(e_0) = 0 \implies \psi(x) = \psi(e_0)\phi(x)\]
  Ceci étant vrai pour tout \(x \in E\), en posant
  \(\lambda = \psi(e_0)\), on obtient l'égalité des formes linéaires
  \(\psi = \lambda \phi\).
\item
  \textbf{Conclusion :} Les deux formes linéaires sont proportionnelles.
\end{enumerate}

\(\blacksquare\) \# Exercice 5: Équation d'un hyperplan en dimension n
\#\# Énoncé Soit \(E\) un espace vectoriel de dimension \(n\) et
\((e_1, \dots, e_n)\) une base de \(E\). Montrer qu'un hyperplan \(H\)
est caractérisé par une équation cartésienne de la forme
\(\sum_{i=1}^n a_i x_i = 0\) où les \((a_i)\) ne sont pas tous nuls.


\subsection{Correction détaillée}

\begin{enumerate}


\item
  \textbf{Caractérisation par une forme linéaire :} Par définition, un
  hyperplan \(H\) est le noyau d'une forme linéaire non nulle, notons-la
  \(\phi : E \to \mathbb{K}\). \[H = \{ x \in E \mid \phi(x) = 0 \}\]
\item
  \textbf{Décomposition dans la base :} Soit un vecteur quelconque
  \(x \in E\). Il admet une unique décomposition dans la base
  \((e_1, \dots, e_n)\) : \[x = \sum_{i=1}^n x_i e_i\] où les \(x_i\)
  sont les coordonnées de \(x\).
\item
  \textbf{Application de la linéarité :} Appliquons \(\phi\) au vecteur
  \(x\) : \[\phi(x) = \phi\left(\sum_{i=1}^n x_i e_i\right)\] Par
  linéarité de \(\phi\), la somme et les scalaires sortent :
  \[\phi(x) = \sum_{i=1}^n x_i \phi(e_i)\]
\item
  \textbf{Identification des coefficients :} Posons pour tout
  \(i \in \{1, \dots, n\}\), \(a_i = \phi(e_i)\). L'équation
  d'appartenance à \(H\) devient alors : \[\sum_{i=1}^n a_i x_i = 0\]
\item
  \textbf{Non-nullité des coefficients :} Comme \(\phi\) est une forme
  linéaire non nulle, il existe au moins un vecteur de base \(e_{i_0}\)
  tel que \(\phi(e_{i_0}) \neq 0\). Donc, il existe au moins un \(a_i\)
  tel que \(a_i \neq 0\).
\item
  \textbf{Conclusion :} Tout hyperplan est rigoureusement caractérisé
  par une équation linéaire homogène dont les coefficients ne sont pas
  tous nuls.
\end{enumerate}

\(\blacksquare\) \# Exercice 6: Bidual d'un espace vectoriel \#\# Énoncé
Pour tout \(x \in E\), on définit l'application d'évaluation
\(\text{ev}_x : E^* \to \mathbb{K}\) par
\(\text{ev}_x(\phi) = \phi(x)\). Montrer que l'application
\(\Psi : x \mapsto \text{ev}_x\) est linéaire et injective.


\subsection{Correction détaillée}

\begin{enumerate}


\item
  \textbf{Étape 1:} \(\text{ev}_x\) est bien une forme linéaire sur
  \(E^*\) car
  \(\text{ev}_x(\phi_1 + \lambda \phi_2) = (\phi_1 + \lambda \phi_2)(x) = \phi_1(x) + \lambda \phi_2(x) = \text{ev}_x(\phi_1) + \lambda \text{ev}_x(\phi_2)\).
\item
  \textbf{Étape 2:} L'application \(\Psi\) est linéaire. Soit
  \(x, y \in E\) et \(\lambda, \mu \in \mathbb{K}\). Évaluons
  \(\Psi(\lambda x + \mu y)\) sur un élément quelconque \(\phi \in E^*\)
  :
  \[\Psi(\lambda x + \mu y)(\phi) = \phi(\lambda x + \mu y) = \lambda \phi(x) + \mu \phi(y) = \lambda \Psi(x)(\phi) + \mu \Psi(y)(\phi)\]
  Ce qui démontre la linéarité.
\item
  \textbf{Étape 3:} Montrons que \(\Psi\) est injective. Soit
  \(x \in \ker \Psi\). Alors pour toute forme linéaire \(\phi \in E^*\),
  on a \(\phi(x) = 0\).
\item
  \textbf{Étape 4:} Si \(x \neq 0\), on pourrait le compléter en une
  base \((x, e_2, \dots, e_n)\) et définir la forme coordonnée
  \(\phi=x^*\) telle que \(x^*(x)=1\), ce qui contredit \(\phi(x)=0\).
\item
  \textbf{Conclusion:} Par conséquent \(x=0\), d'où
  \(\ker \Psi = \{0\}\). \(\Psi\) est injective.
\end{enumerate}

\(\blacksquare\) \# Exercice 7: Trace comme forme linéaire \#\# Énoncé
L'application \(\text{Tr} : M_n(\mathbb{K}) \to \mathbb{K}\) est une
forme linéaire. Montrer que tout hyperplan \(H\) de \(M_n(\mathbb{K})\)
contient au moins une matrice inversible (pour \(n \ge 2\)).


\subsection{Correction détaillée}

\begin{enumerate}


\item
  \textbf{Étape 1:} Un hyperplan de \(M_n(\mathbb{K})\) est le noyau
  d'une forme linéaire non nulle \(\phi\). Il est connu que toute forme
  linéaire sur \(M_n(\mathbb{K})\) s'écrit \(\phi(M) = \text{Tr}(AM)\)
  pour une unique matrice \(A \in M_n(\mathbb{K})\). Ainsi,
  \(H = \{ M \in M_n(\mathbb{K}) \mid \text{Tr}(AM) = 0 \}\).
\item
  \textbf{Étape 2:} On cherche \(M \in H\) telle que \(\det(M) \neq 0\).
  Si \(A=0\), \(\phi=0\) ce qui est exclu. Supposons par l'absurde que
  \(H\) ne contient aucune matrice inversible.
\item
  \textbf{Étape 3:} L'hyperplan \(H\) est un sous-espace vectoriel de
  dimension \(n^2-1\). Si \(H\) ne contient que des matrices
  singulières, on a une contradiction avec les résultats de la théorie
  des espaces de matrices de rang borné (théorème de Dieudonné), car la
  dimension maximale d'un sous-espace de matrices non-inversibles est
  \(n(n-1)\), et pour \(n \ge 2\), \(n^2-1 > n(n-1)\).
\item
  \textbf{Conclusion:} L'hypothèse de départ est fausse, donc on en
  déduit formellement qu'un hyperplan contient toujours des éléments
  inversibles.
\end{enumerate}

\(\blacksquare\) \# Exercice 8: Orthogonalité croisée \#\# Énoncé Soient
\(F_1, F_2\) deux sous-espaces de \(E\). Montrer que
\((F_1 + F_2)^\perp = F_1^\perp \cap F_2^\perp\).


\subsection{Correction détaillée}

\begin{enumerate}


\item
  \textbf{Étape 1 :} Soit \(\phi \in (F_1 + F_2)^\perp\). Alors pour
  tout \(x \in F_1 + F_2, \phi(x)=0\). En particulier, pour tout
  \(x_1 \in F_1\), \(\phi(x_1)=0\) (car \(F_1 \subset F_1+F_2\)), donc
  \(\phi \in F_1^\perp\). De même \(\phi \in F_2^\perp\). Ainsi
  \((F_1 + F_2)^\perp \subset F_1^\perp \cap F_2^\perp\).
\item
  \textbf{Étape 2 :} Réciproquement, soit
  \(\phi \in F_1^\perp \cap F_2^\perp\). Pour tout \(x \in F_1+F_2\), on
  peut écrire \(x = x_1 + x_2\) avec \(x_1 \in F_1, x_2 \in F_2\).
\item
  \textbf{Étape 3 :} On a par linéarité
  \(\phi(x) = \phi(x_1+x_2) = \phi(x_1) + \phi(x_2)\).
\item
  \textbf{Étape 4 :} Puisque \(\phi \in F_1^\perp\), \(\phi(x_1)=0\).
  Puisque \(\phi \in F_2^\perp\), \(\phi(x_2)=0\). Donc
  \(\phi(x) = 0 + 0 = 0\). Ainsi \(\phi \in (F_1+F_2)^\perp\). On a
  l'inclusion réciproque.
\item
  \textbf{Conclusion:} Par double inclusion, l'égalité
  \((F_1 + F_2)^\perp = F_1^\perp \cap F_2^\perp\) est strictement
  démontrée.
\end{enumerate}

\(\blacksquare\) \# Exercice 9: Indépendance linéaire de formes
linéaires \#\# Énoncé Soient \(\phi_1, \dots, \phi_p \in E^*\). Montrer
qu'elles sont linéairement indépendantes si et seulement si
l'intersection de leurs noyaux \(\bigcap_{i=1}^p \ker \phi_i\) est de
dimension \(n-p\).


\subsection{Correction détaillée}

\begin{enumerate}


\item
  \textbf{Étape 1:} On définit l'application linéaire
  \(\Phi : E \to \mathbb{K}^p\) par
  \(\Phi(x) = (\phi_1(x), \dots, \phi_p(x))\).
\item
  \textbf{Étape 2:} Le noyau de \(\Phi\) est exactement l'intersection
  des noyaux : \(\ker \Phi = \bigcap_{i=1}^p \ker \phi_i\).
\item
  \textbf{Étape 3:} D'après le théorème du rang,
  \(\dim E = \dim \ker \Phi + \dim \text{Im}(\Phi)\). Donc
  \(\dim \bigcap \ker \phi_i = n - \dim \text{Im}(\Phi)\).
\item
  \textbf{Étape 4:} L'image de la transposée
  \(\Phi^t : (\mathbb{K}^p)^* \to E^*\) est le sous-espace engendré par
  les \(\phi_i\). Or
  \(\text{rg}(\Phi) = \text{rg}(\Phi^t) = \dim \text{Vect}(\phi_1, \dots, \phi_p)\).
\item
  \textbf{Étape 5:} Ainsi,
  \(\dim \text{Vect}(\phi_1, \dots, \phi_p) = p\) (c'est-à-dire que la
  famille est libre) si et seulement si \(\text{rg}(\Phi) = p\), ce qui
  équivaut à \(\dim \bigcap \ker \phi_i = n - p\).
\item
  \textbf{Conclusion:} L'équivalence est rigoureusement prouvée via
  l'application associée et le théorème du rang.
\end{enumerate}

\(\blacksquare\) \# Exercice 10: Polynômes de Lagrange et dualité \#\#
Énoncé Dans \(E = \mathbb{R}_{n-1}[X]\), on se donne \(n\) scalaires
distincts \(a_1, \dots, a_n\). Montrer que les formes linéaires
d'évaluation \(\phi_i(P) = P(a_i)\) forment une base de \(E^*\).


\subsection{Correction détaillée}

\begin{enumerate}


\item
  \textbf{Étape 1:} La dimension de \(E\) est \(n\), donc
  \(\dim E^* = n\). Pour démontrer qu'une famille de \(n\) vecteurs
  forme une base, il suffit de montrer que la famille
  \((\phi_1, \dots, \phi_n)\) est libre.
\item
  \textbf{Étape 2:} Supposons une combinaison linéaire nulle:
  \(\sum_{i=1}^n \lambda_i \phi_i = 0\). Cela signifie que pour tout
  polynôme \(P \in E\), \(\sum_{i=1}^n \lambda_i P(a_i) = 0\).
\item
  \textbf{Étape 3:} On introduit les polynômes interpolateurs de
  Lagrange \(L_j(X) = \prod_{k \neq j} \frac{X-a_k}{a_j-a_k}\). Ce sont
  des éléments de \(E\) car leur degré est exactement \(n-1\).
\item
  \textbf{Étape 4:} Par construction, \(L_j(a_i) = \delta_{ij}\) (vaut 1
  si \(i=j\), 0 sinon).
\item
  \textbf{Étape 5:} Évaluons la combinaison linéaire nulle sur le
  polynôme \(L_j\) :
  \[0 = \sum_{i=1}^n \lambda_i \phi_i(L_j) = \sum_{i=1}^n \lambda_i L_j(a_i) = \lambda_j\]
\item
  \textbf{Conclusion:} Pour tout indice \(j\), on obtient
  \(\lambda_j = 0\). La famille est donc libre, et c'est par conséquent
  une base. Les polynômes \((L_j)\) forment précisément la base
  antéduale associée.
\end{enumerate}

\(\blacksquare\)


\section{Travaux Pratiques et Simulations
Algorithmiques}


\section{TP 1: Implémentation pure Python d'une forme
linéaire}


\subsection{Objectif mathématique}

Mise en pratique algorithmique de l'espace dual, de l'hyperplan et des
notions fondamentales de l'algèbre linéaire, structurée en Python absolu
sans bibliothèques de haut niveau.


\subsection{Code Python}

\begin{Shaded}
\begin{Highlighting}[]
\KeywordTok{def}\NormalTok{ apply\_linear\_form(form\_coeffs, vector):}
    \ControlFlowTok{if} \BuiltInTok{len}\NormalTok{(form\_coeffs) }\OperatorTok{!=} \BuiltInTok{len}\NormalTok{(vector):}
        \ControlFlowTok{raise} \PreprocessorTok{ValueError}\NormalTok{(}\StringTok{"Dimensions mismatch"}\NormalTok{)}

    \CommentTok{\# Implémentation mathématique de la somme des produits (produit scalaire)}
\NormalTok{    resultat }\OperatorTok{=} \DecValTok{0}
    \ControlFlowTok{for}\NormalTok{ i }\KeywordTok{in} \BuiltInTok{range}\NormalTok{(}\BuiltInTok{len}\NormalTok{(vector)):}
\NormalTok{        resultat }\OperatorTok{+=}\NormalTok{ form\_coeffs[i] }\OperatorTok{*}\NormalTok{ vector[i]}
    \ControlFlowTok{return}\NormalTok{ resultat}

\CommentTok{\# Validation avec un calcul manuel: 2(1) {-} 1(2) + 3(0) = 0}
\ControlFlowTok{assert}\NormalTok{ apply\_linear\_form([}\DecValTok{2}\NormalTok{, }\OperatorTok{{-}}\DecValTok{1}\NormalTok{, }\DecValTok{3}\NormalTok{], [}\DecValTok{1}\NormalTok{, }\DecValTok{2}\NormalTok{, }\DecValTok{0}\NormalTok{]) }\OperatorTok{==} \DecValTok{0}
\ControlFlowTok{assert}\NormalTok{ apply\_linear\_form([}\DecValTok{1}\NormalTok{, }\DecValTok{1}\NormalTok{], [}\DecValTok{3}\NormalTok{, }\DecValTok{4}\NormalTok{]) }\OperatorTok{==} \DecValTok{7}
\end{Highlighting}
\end{Shaded}


\subsection{Validation Rigoureuse}

Ce code vérifie le calcul algorithmique fondamental d'une forme linéaire
\(\phi(x) = \sum a_i x_i\) par une boucle itérative stricte, sans
artifice de bibliothèque externe. \# TP 2: Recherche de l'équation d'un
hyperplan affine


\subsection{Objectif mathématique}

Mise en pratique algorithmique de l'espace dual, de l'hyperplan et des
notions fondamentales de l'algèbre linéaire, structurée en Python absolu
sans bibliothèques de haut niveau.


\subsection{Code Python}

\begin{Shaded}
\begin{Highlighting}[]
\KeywordTok{def}\NormalTok{ hyperplane\_equation\_from\_normal(normal\_vector, point\_on\_plane):}
    \CommentTok{\# L\textquotesingle{}équation de l\textquotesingle{}hyperplan est sum(n\_i * x\_i) = d}
    \CommentTok{\# Pour trouver d, on évalue la forme linéaire sur le point appartenant au plan.}
\NormalTok{    d }\OperatorTok{=} \DecValTok{0}
    \ControlFlowTok{for}\NormalTok{ i }\KeywordTok{in} \BuiltInTok{range}\NormalTok{(}\BuiltInTok{len}\NormalTok{(normal\_vector)):}
\NormalTok{        d }\OperatorTok{+=}\NormalTok{ normal\_vector[i] }\OperatorTok{*}\NormalTok{ point\_on\_plane[i]}
    \ControlFlowTok{return}\NormalTok{ normal\_vector, d}

\NormalTok{normal, constante\_d }\OperatorTok{=}\NormalTok{ hyperplane\_equation\_from\_normal([}\DecValTok{1}\NormalTok{, }\OperatorTok{{-}}\DecValTok{1}\NormalTok{, }\DecValTok{2}\NormalTok{], [}\DecValTok{1}\NormalTok{, }\DecValTok{1}\NormalTok{, }\DecValTok{1}\NormalTok{])}
\CommentTok{\# Calcul formel: 1(1) {-} 1(1) + 2(1) = 2}
\ControlFlowTok{assert}\NormalTok{ normal }\OperatorTok{==}\NormalTok{ [}\DecValTok{1}\NormalTok{, }\OperatorTok{{-}}\DecValTok{1}\NormalTok{, }\DecValTok{2}\NormalTok{]}
\ControlFlowTok{assert}\NormalTok{ constante\_d }\OperatorTok{==} \DecValTok{2}
\end{Highlighting}
\end{Shaded}


\subsection{Validation Rigoureuse}

Ce code calcule rigoureusement la constante affine \(d\) de l'hyperplan
en traduisant l'égalité \(\phi(x_0) = d\) en Python pur. \# TP 3:
Génération algorithmique d'une base de noyau (Dimension 3)


\subsection{Objectif mathématique}

Mise en pratique algorithmique de l'espace dual, de l'hyperplan et des
notions fondamentales de l'algèbre linéaire, structurée en Python absolu
sans bibliothèques de haut niveau.


\subsection{Code Python}

\begin{Shaded}
\begin{Highlighting}[]
\KeywordTok{def}\NormalTok{ find\_hyperplane\_basis\_3d(coeffs):}
\NormalTok{    a, b, c }\OperatorTok{=}\NormalTok{ coeffs}
    \CommentTok{\# Recherche systématique d\textquotesingle{}une base libre}
    \ControlFlowTok{if}\NormalTok{ a }\OperatorTok{!=} \DecValTok{0}\NormalTok{:}
        \CommentTok{\# On exprime x en fonction de y et z: x = ({-}b/a)y + ({-}c/a)z}
        \ControlFlowTok{return}\NormalTok{ [(}\OperatorTok{{-}}\NormalTok{b}\OperatorTok{/}\NormalTok{a, }\DecValTok{1}\NormalTok{, }\DecValTok{0}\NormalTok{), (}\OperatorTok{{-}}\NormalTok{c}\OperatorTok{/}\NormalTok{a, }\DecValTok{0}\NormalTok{, }\DecValTok{1}\NormalTok{)]}
    \ControlFlowTok{elif}\NormalTok{ b }\OperatorTok{!=} \DecValTok{0}\NormalTok{:}
        \ControlFlowTok{return}\NormalTok{ [(}\DecValTok{1}\NormalTok{, }\OperatorTok{{-}}\NormalTok{a}\OperatorTok{/}\NormalTok{b, }\DecValTok{0}\NormalTok{), (}\DecValTok{0}\NormalTok{, }\OperatorTok{{-}}\NormalTok{c}\OperatorTok{/}\NormalTok{b, }\DecValTok{1}\NormalTok{)]}
    \ControlFlowTok{elif}\NormalTok{ c }\OperatorTok{!=} \DecValTok{0}\NormalTok{:}
        \ControlFlowTok{return}\NormalTok{ [(}\DecValTok{1}\NormalTok{, }\DecValTok{0}\NormalTok{, }\OperatorTok{{-}}\NormalTok{a}\OperatorTok{/}\NormalTok{c), (}\DecValTok{0}\NormalTok{, }\DecValTok{1}\NormalTok{, }\OperatorTok{{-}}\NormalTok{b}\OperatorTok{/}\NormalTok{c)]}
    \ControlFlowTok{else}\NormalTok{:}
        \ControlFlowTok{raise} \PreprocessorTok{ValueError}\NormalTok{(}\StringTok{"La forme linéaire ne peut être identiquement nulle"}\NormalTok{)}

\NormalTok{basis }\OperatorTok{=}\NormalTok{ find\_hyperplane\_basis\_3d([}\DecValTok{2}\NormalTok{, }\OperatorTok{{-}}\DecValTok{1}\NormalTok{, }\DecValTok{3}\NormalTok{])}
\CommentTok{\# Validation mathématique: la forme linéaire s\textquotesingle{}annule bien sur les vecteurs de base}
\ControlFlowTok{assert} \BuiltInTok{sum}\NormalTok{(c}\OperatorTok{*}\NormalTok{v }\ControlFlowTok{for}\NormalTok{ c, v }\KeywordTok{in} \BuiltInTok{zip}\NormalTok{([}\DecValTok{2}\NormalTok{, }\OperatorTok{{-}}\DecValTok{1}\NormalTok{, }\DecValTok{3}\NormalTok{], basis[}\DecValTok{0}\NormalTok{])) }\OperatorTok{==} \DecValTok{0}
\ControlFlowTok{assert} \BuiltInTok{sum}\NormalTok{(c}\OperatorTok{*}\NormalTok{v }\ControlFlowTok{for}\NormalTok{ c, v }\KeywordTok{in} \BuiltInTok{zip}\NormalTok{([}\DecValTok{2}\NormalTok{, }\OperatorTok{{-}}\DecValTok{1}\NormalTok{, }\DecValTok{3}\NormalTok{], basis[}\DecValTok{1}\NormalTok{])) }\OperatorTok{==} \DecValTok{0}
\end{Highlighting}
\end{Shaded}


\subsection{Validation Rigoureuse}

La fonction construit explicitement une famille génératrice libre pour
le noyau (sous-espace de dimension \(n-1\)), validée par assertion. \#
TP 4: Intersection analytique de deux hyperplans dans R\^{}2


\subsection{Objectif mathématique}

Mise en pratique algorithmique de l'espace dual, de l'hyperplan et des
notions fondamentales de l'algèbre linéaire, structurée en Python absolu
sans bibliothèques de haut niveau.


\subsection{Code Python}

\begin{Shaded}
\begin{Highlighting}[]
\KeywordTok{def}\NormalTok{ intersect\_2d\_lines(a1, b1, d1, a2, b2, d2):}
    \CommentTok{\# Calcul du déterminant du système}
\NormalTok{    det }\OperatorTok{=}\NormalTok{ a1b1}
    \ControlFlowTok{if}\NormalTok{ det }\OperatorTok{==} \DecValTok{0}\NormalTok{:}
        \ControlFlowTok{return} \VariableTok{None} \CommentTok{\# Les hyperplans (droites) sont parallèles ou confondus}

    \CommentTok{\# Résolution exacte par la règle de Cramer}
\NormalTok{    x }\OperatorTok{=}\NormalTok{ (d1b1) }\OperatorTok{/}\NormalTok{ det}
\NormalTok{    y }\OperatorTok{=}\NormalTok{ (a1d1) }\OperatorTok{/}\NormalTok{ det}
    \ControlFlowTok{return}\NormalTok{ (x, y)}

\CommentTok{\# Intersection de x {-} y = 0 et x + y = 2}
\ControlFlowTok{assert}\NormalTok{ intersect\_2d\_lines(}\DecValTok{1}\NormalTok{, }\OperatorTok{{-}}\DecValTok{1}\NormalTok{, }\DecValTok{0}\NormalTok{, }\DecValTok{1}\NormalTok{, }\DecValTok{1}\NormalTok{, }\DecValTok{2}\NormalTok{) }\OperatorTok{==}\NormalTok{ (}\FloatTok{1.0}\NormalTok{, }\FloatTok{1.0}\NormalTok{)}
\end{Highlighting}
\end{Shaded}


\subsection{Validation Rigoureuse}

L'intersection de deux hyperplans dans \(\mathbb{R}^2\) est résolue de
manière algorithmique absolue par la formule de Cramer. \# TP 5: Dualité
: Translation Vecteur-Forme (Théorème de Riesz discret)


\subsection{Objectif mathématique}

Mise en pratique algorithmique de l'espace dual, de l'hyperplan et des
notions fondamentales de l'algèbre linéaire, structurée en Python absolu
sans bibliothèques de haut niveau.


\subsection{Code Python}

\begin{Shaded}
\begin{Highlighting}[]
\KeywordTok{def}\NormalTok{ riesz\_representation(form\_eval\_func, dimension):}
    \CommentTok{\# La forme linéaire est entièrement déterminée par ses valeurs sur la base canonique.}
    \CommentTok{\# On reconstruit le vecteur représentant \textasciigrave{}u\textasciigrave{} tel que phi(x) = \textless{}u, x\textgreater{}}

\NormalTok{    representing\_vector }\OperatorTok{=}\NormalTok{ []}
    \ControlFlowTok{for}\NormalTok{ i }\KeywordTok{in} \BuiltInTok{range}\NormalTok{(dimension):}
        \CommentTok{\# Création du i{-}ème vecteur de la base canonique e\_i}
\NormalTok{        e\_i }\OperatorTok{=}\NormalTok{ [}\DecValTok{0}\NormalTok{] }\OperatorTok{*}\NormalTok{ dimension}
\NormalTok{        e\_i[i] }\OperatorTok{=} \DecValTok{1}

        \CommentTok{\# Le coefficient a\_i est l\textquotesingle{}évaluation de la forme sur e\_i}
\NormalTok{        a\_i }\OperatorTok{=}\NormalTok{ form\_eval\_func(e\_i)}
\NormalTok{        representing\_vector.append(a\_i)}

    \ControlFlowTok{return}\NormalTok{ representing\_vector}

\KeywordTok{def}\NormalTok{ custom\_form(v):}
    \CommentTok{\# Une forme linéaire quelconque agissant sur un vecteur 3D}
    \ControlFlowTok{return} \DecValTok{3}\ErrorTok{v}\NormalTok{[}\DecValTok{1}\NormalTok{] }\OperatorTok{+} \DecValTok{5}\OperatorTok{*}\NormalTok{v[}\DecValTok{2}\NormalTok{]}

\NormalTok{repr\_vec }\OperatorTok{=}\NormalTok{ riesz\_representation(custom\_form, }\DecValTok{3}\NormalTok{)}
\ControlFlowTok{assert}\NormalTok{ repr\_vec }\OperatorTok{==}\NormalTok{ [}\DecValTok{3}\NormalTok{, }\OperatorTok{{-}}\DecValTok{2}\NormalTok{, }\DecValTok{5}\NormalTok{]}
\end{Highlighting}
\end{Shaded}


\subsection{Validation Rigoureuse}

Démontre en Python pur la construction de l'isomorphisme de Riesz :
extraction rigoureuse du vecteur représentatif en évaluant la forme sur
la base canonique.

\vspace{1cm}
\begin{center}
\begin{tikzpicture}[scale=1.5]
    % Axes
    \draw[->] (-2,0) -- (3,0) node[right] {$x$};
    \draw[->] (0,-1) -- (0,3) node[above] {$y$};

    % Hyperplane (Line in 2D)
    \draw[thick, blue] (-1.5, -0.5) -- (2.5, 2.166) node[above right] {$\ker(\phi)$ (Hyperplan $H$)};

    % Normal vector
    \draw[->, thick, red] (0.5, 0.833) -- (0.0, 1.583) node[above left] {Vecteur normal $w$};

    % Point and projection
    \filldraw[black] (2, 0.5) circle (1pt) node[below right] {$v$};

    \node[align=center, below] at (0.5, -1.5) {Illustration de la dualité et séparation spatiale};
\end{tikzpicture}
\end{center}

\end{document}
